% ============================================================================
%  THE SYNTHESIS PAPER — ARC 18.  Written from SYNTHESIS_HARVEST.md and from
%  nothing else: every claim below traces to a harvested entry with its home,
%  its anchor, and the register tag its source gave it.
%
%  THREE RULES THIS FILE IS WRITTEN UNDER, and they are the arc's own:
%    * ONE STATE.  No "was open, later resolved" anywhere.  The fact of a
%      resolution belongs in the consolidation record, not in a results paper.
%    * EVERY "IS" AUDITED.  A claim is stated at the weight its source states
%      it: a proposition as a proposition, a reading as a reading, a
%      conjecture as a conjecture.  The harvest's register field is the
%      authority for that, not this paper's convenience.
%    * THE NEGATIVES TRAVEL.  Where a harvested entry carries a bound in its
%      "not operative for" field, the bound is stated beside the claim.  A
%      distillation that drops those has not shortened the corpus; it has
%      misreported it.
% ============================================================================
\documentclass[11pt,a4paper]{article}
\usepackage[margin=1in]{geometry}
\usepackage{amsmath, amssymb, amsthm}
\usepackage{mathtools}
\usepackage{graphicx}
\usepackage[dvipsnames]{xcolor}
\usepackage{hyperref}
\usepackage{caption}
\usepackage{receipts}
\usepackage{ledgers}

\newcommand{\dd}{\mathrm{d}}
\newcommand{\dS}{\mathrm{dS}}
\newcommand{\so}{\mathfrak{so}}
\newcommand{\rs}{r_{s}}

\theoremstyle{plain}
\newtheorem{proposition}{Proposition}
\newtheorem{theorem}[proposition]{Theorem}
\theoremstyle{definition}
\newtheorem{remark}[proposition]{Remark}

\title{Cosmological Relativity: what the construction delivers,\\
what it declines, and where its edge is}
\author{D.~Janzen}
\date{}

\begin{document}
\maketitle

\begin{abstract}
This paper states, in one place, what the Cosmological Relativity corpus establishes, at the weight each
result is established at, together with what the construction declines to claim and where its open edge
lies. Nothing here is new: every claim is drawn from one of seventeen papers or from the ledgers those
papers cite, and is stated at the register its source gives it---a proposition as a proposition, a reading
as a reading, a conjecture as a conjecture. What is new is the assembly, and the two things an assembly can
supply that a set of papers cannot: the joins, where a result in one paper is load-bearing for a result in
another, and the relative size of the frontier against the body of work that surrounds it.
\end{abstract}

\section{What forces the object, what it is, and the distinctions it forces}\label{sec:object}

Cosmological Relativity is a construction on one geometric object, and almost everything below is that
object read from somewhere. But the object is not where the argument starts, and presenting it first would
ask a reader to grant the thing in question. It is reached: a measurement forces a cosmic foliation, reading
that foliation ontologically is the augmentation, and maximal symmetry then leaves exactly one manifold to
carry it. This section takes those in order, and closes on the distinction every later result depends on ---
what a \emph{boundary} of such an object is.

\subsection{What the measurement forces}

The cosmological redshift is not a property an emitting source carries. It is the path integral of the
expansion rate along the photon's own trajectory, accumulated continuously rather than fixed at
emission~\cite{JanzenModernParallax}. That distinction is what makes the isotropy of the microwave
background's monopole a measurement of something other than what it is usually taken to measure.

\emph{Grant the common reading its strongest form.} Let last scattering be \emph{perfectly} homogeneous.
Then the observed temperature contrast is exactly the variation in the integrated expansion, with nothing
left over, and \emph{the entire observed isotropy is the statement that the integrated expansion was the
same in every direction}. Homogeneity at the source does not \emph{supply} the isotropy; granted perfectly,
it removes one term and leaves the isotropy as a direct measurement of uniform
expansion~\cite{JanzenModernParallax}.

Under the alternative --- that the large-scale expansion rate tracks the matter lumpiness region by region,
with no single global scale factor --- the accumulated redshift would scatter across the sky by
$\sim10^{-3}$, against an observed $\lesssim3\times10^{-6}$: \emph{an exclusion by three orders of
magnitude}. \emph{And every choice in the estimate biases it downward, so the number is a floor rather than
an estimate}~\cite{JanzenModernParallax}.

Two escapes remain and they are closed separately. The \emph{statistical} escape is closed by the floor. The
\emph{structured} escape --- a finely tuned, spherically symmetric inhomogeneity centred on the observer,
which would give an isotropic redshift with no global uniform expansion and is untouched by any
$\sqrt N$ argument --- is closed by the Copernican principle \emph{together with} the independently measured
isotropy of the expansion history itself. \emph{With both excluded, the observed isotropy forces uniform
cosmic expansion.}

\emph{What a single datum thereby establishes is worth stating exactly, because it is larger than a
parameter bound.} The rejected region is not a corner of parameter space but \emph{the entire class of
histories lacking a single global scale factor on a single global time}. To reject that class is to select a
global, time-ordered expansion --- the cosmic foliation, made dynamical. \emph{And the foliation is
logically prior to everything the standard model otherwise assumes}: ``space'' is one of its slices,
``expansion'' an ordered family of them, isotropy and homogeneity properties of a slice --- none of them
statable until the foliation is fixed.

\subsection{The augmentation that follows, and its structural counterpart}

Augmenting general relativity by fixing a physical foliation and reading it ontologically --- the lapse the
objective rate at which the existing world advances, the shift the relativity of synchrony --- is
\emph{necessary and sufficient} for a coherent formal description of an existing, evolving
world~\cite{JanzenCRframework}. \emph{Necessary}, because a description lacking the lapse part collapses
existence into occurrence. \emph{Sufficient}, because the two parts close the only coherent escapes. No field
equation is altered.

\emph{And the same augmentation is forced a second time, from general relativity alone.} The event horizon
of a collapsing body is a metric singularity met only at infinite exterior time, so no finite layer carries
a completed horizon (\S\ref{sec:boundary})~\cite{JanzenBHcausality}. \emph{These two forcings are
independent and co-equal, and neither rests on the other}: one is measured and one is structural, and a
reader may take either without the other.

\subsection{Three senses of ``real'', fixed before the word is used}

The corpus uses \emph{real} in three senses and keeps them apart; this paper does the same. They are not
degrees of one thing.

\emph{Real in the arithmetic sense}: a geometry has a real coordinate basis, with no imaginary coordinate
doing any work. That is the sense in which the substrate below is a real manifold, and the claim is about
coordinates and signature and nothing more.

\emph{Real by construction}: a feature is built by a construction and really seen from the vantage that
builds it. The perspectival mass is of this kind, and so is the compact face on which a continuous
$\mathfrak{su}(3)$ could live.

\emph{Existent}: the ontological sense, whose criterion in this framework is temporality. \emph{The existent
is the evolving three-dimensional layer, and the manifold is its record}~\cite{JanzenCRframework}. By that
criterion the compact face is real by construction but, being Riemannian and atemporal, \emph{not a co-equal
existent}---it holds no clock~\cite{JanzenBoundary}.

\emph{So nothing in this section claims that a five-dimensional object is what exists.} The substrate is the
geometric core the construction reads and cuts; what exists, on this framework's own criterion, is the layer.
The higher-dimensional geometry is what the layer is a cut of, not a thing possessing the layer's kind of
being, and the four-dimensional manifold is the record rather than the existent.

\subsection{What maximal symmetry then selects}

\emph{The foliation is forced and the manifold that carries it is not yet fixed, and this is where the
construction's one object arrives --- as a selection rather than a premise.} The maximally symmetric
solutions of $R_{\mu\nu}=\Lambda g_{\mu\nu}$ form a one-parameter family with curvature radius
$\alpha=\sqrt{3/|\Lambda|}$. They are usually met as hypersurfaces in a flat five-dimensional space, which is
a convenient picture and not the manifold: solving that embedding for the fifth coordinate and substituting
back returns an intrinsic four-dimensional line element in which the manifold is coordinatised by four
\emph{real} coordinates, real regardless of the sign of the curvature.

The imaginary fifth coordinate is a property of one embedding picture, not of the surface. \emph{This is reality in the first of the
three senses above, and in that one only.}

Reading the metric's eigenvalues off that real line element gives the result the programme's ontology stands
on. Three eigenvalues are positive always, and the fourth's sign is fixed by real data alone.

\begin{proposition}[The unique intrinsically Lorentzian maximally symmetric manifold]\label{prop:unique}
De~Sitter space is the only real maximally symmetric manifold carrying an intrinsic Lorentzian signature.
It has a real coordinate basis, satisfies the maximal-symmetry requirement, and is the only such manifold
with the requisite causal structure~\cite{JanzenGeometricCore,JanzenThesis}.
\end{proposition}

Two ways the signature can turn are worth separating, because they are of different kinds. At fixed positive
curvature the crossing occurs where the fourth eigenvalue passes through a \emph{pole}: its magnitude
diverges and its sign flips, the numerator being constant throughout. \emph{So the metric does not degenerate
at the crossing}---a signature change through zero would send the determinant to zero and leave no metric
there, and that is not what happens. The one place the eigenvalue vanishes is the null cone: not a crossing
within a member of the family, but the family's own singular leaf~\cite{JanzenGeometricCore}.

That the substrate is reached through imaginary instruments is therefore a fact about the instruments. The
embedding coordinate, the equatorial seam's continuation, and the reassignment at the cosmogenesis are each
places the construction uses an imaginary variable and lands on a real manifold; the geometry is
intrinsically real at every point they land on~\cite{JanzenGeometricCore}. In one case the instrument is not
optional and the case is the construction's own central algebraic object: over the field of the mass
parameter the horizon cubic is irreducible, and below the Nariai mass its discriminant is positive, so its
three roots are real and distinct. That is exactly the hypothesis pair of \emph{casus irreducibilis}, whose
conclusion is that no root lies in any real radical extension---\emph{the three real horizon radii admit no
expression in real radicals as functions of the mass}, and every radical route to them passes through
$\mathbb{C}$~\cite{JanzenGeometricCore}. The imaginary is the only available road, and the roots it lands on
are real.

\subsection{One scale, and what that means}

The substrate carries a single dimensionful scale. Read projectively it is a Cayley--Klein
geometry\ldg{quadric_geometry} whose \emph{absolute} is the null cone, and its geodesic separation is the
Cayley--Klein log-cross-ratio taken against that absolute, with the Cayley--Klein constant equal to
$\alpha$~\cite{JanzenGeometricCore}. So $\alpha$ is not a length inserted into a metric but the constant of a
projective geometry whose absolute is the light cone---which is the signature proposition above, read
projectively rather than metrically.

The constants through which physics is written on it are unit gauges, each a nameable feature of the
substrate: $c$ is the null-ruling slope of the equilateral hyperboloid, $\Lambda$ the waist, and $G$ enters
only as the offset-length $GM/c^{2}$, the mass being the offset of the cut from the central
geodesic~\cite{JanzenGeometricCore,JanzenOperator,JanzenSlicing}. On this accounting the
gravitational--cosmological--quantum sector spends no free dimensionless constant, and the reason is
structural rather than counted: neither real form supplies a second invariant, and a dimensionless magnitude
needs two~\cite{JanzenGeometricCore}.

\emph{What the count leaves is not a constant but an initial condition}, and the corpus carries it as
measured content rather than as a derived quantity---as flat $\Lambda$CDM carries the baryon-to-photon ratio
from a baryogenesis it does not model~\cite{JanzenCRcosmology}.

\subsection{What a boundary of the substrate is}\label{sec:boundary}

Here the corpus draws the distinction that everything downstream depends on, and it is a distinction between
two kinds of singular locus rather than between a real one and an artefact.

A \emph{metric} singularity is a locus at which a spatial extent contracts to zero while the curvature stays
finite. A curvature singularity is a locus at which an invariant of the curvature diverges. The event
horizon of a collapsing body is the first kind: a null hypersurface along whose generators the spatial extent
also contracts to zero, occurring only as the null future boundary approached at infinite exterior
time~\cite{JanzenBHcausality}. Three consequences of that follow immediately, and each is a distinction a
reader needs before meeting any application.

\emph{Observer-dependence is not the criterion.} It is tempting to sort these loci by whether they are seen
by every observer, and the sorting fails at once: the Rindler horizon is perspectival and carries a flux, so
``perspectival, therefore no flux'' is refuted before it starts. What sorts the four horizons the corpus
treats is \emph{completion}---whether the boundary is reached on any finite layer---and completion sorts all
four~\cite{JanzenBHcausality}.

\emph{And the finite-curvature case is the one the metric-singularity result treats.} It is a different
object from the locus standardly called the curvature singularity at $r=0$, where the invariants computed on
the areal radius diverge; the first is established with the event horizon as its exemplar, and nothing is
said there about the second beyond noting that it is a different object~\cite{JanzenBHcausality}. \emph{That
the two admit a common treatment, and what becomes of the divergence at $r=0$ under it, is established in the
companion}~\cite{JanzenCircle,JanzenSlicing}. \emph{No claim of necessity or sufficiency is made here in
either direction}: the scope is what is stated, and reading it as a general criterion for when a boundary may
be crossed would assert what the source declines to assert.

\emph{And formal likeness does not sort the two classes.} A chart such as Eddington--Finkelstein parametrises
approaches to a singular locus by a coordinate that degenerates there, so a single metric place is drawn
spread out as an extended line. Mercator draws the North Pole by the same projective mechanism, and the
visual appearance is identical: a place drawn as a line. \emph{But the reason for the appearance is the
opposite.} At the pole the chart manufactures an extension that is not there, and a chart centred on the pole
collapses the line back to the point it always was. At the horizon the chart spreads out a metric collapse
that \emph{is} there: reading the line back yields no ordinary point a better chart would reveal, because no
chart removes it---\emph{the collapse is in the metric, not in the projection}~\cite{JanzenCircle}. So no
argument from how a singular locus is drawn can sort the two, in either direction.

\subsection{What follows for the reading, and what does not}

Two guards belong here rather than later, because both are misread in the same direction.

\emph{Perspectival does not mean unreal.} Where the corpus calls a feature perspectival it means built by a
construction and really seen from the vantage that builds it---the curvature divergence of a perspectival
metric is real as that metric's, and reading it as an artefact in the pejorative sense is the misreading the
epistemology is written to defeat~\cite{JanzenShadowExistence,JanzenCircle}.

\emph{And identifying a feature as perspectival is not a dismissal of the construction that carries it.}
Sweeping a radial curve to recover a spatial geometry is the legitimate and indeed the natural way to chart a
geometry from a given vantage, and the result agrees with the standard geometry on every observable outside
the horizon. What is identified is not an illegitimate operation but a misreading of its product: treating
features the construction manufactures---the asymmetric labelling, and with it the curvature singularity at
$r=0$---as features of the geometry rather than of the chart. \emph{The construction is faithful as a
description from its vantage; the error is ontological}~\cite{JanzenCircle}.

\begin{thebibliography}{99}
\bibitem{JanzenThesis} D.~Janzen, \emph{A Solution to the Cosmological Problem of Relativity Theory} (PhD dissertation, University of Saskatchewan, 2012), ch.~3.
\bibitem{JanzenBHcausality} D.~Janzen, \emph{The causal structure of gravitational collapse: the event horizon as a metric singularity}.
\bibitem{JanzenCircle} D.~Janzen, \emph{The Schwarzschild circle: one cycloid, two metric singularities}.
\bibitem{JanzenSlicing} D.~Janzen, \emph{The Schwarzschild--de~Sitter slicing curve}.
\bibitem{JanzenShadowExistence} D.~Janzen, \emph{The shadow of existence: inferring the world that casts the appearances}.
\bibitem{JanzenOperator} D.~Janzen, \emph{The slicing operator: straight cuts are vacuum, matter is the bend of the cut}.
\bibitem{JanzenCRcosmology} D.~Janzen, \emph{The cosmology of Cosmological Relativity and its scalar perturbation sector}.
\bibitem{JanzenModernParallax} D.~Janzen, \emph{A modern parallax: the redshift-isotropy floor and the augmentation it forces}.
\bibitem{JanzenCRframework} D.~Janzen, \emph{Cosmological Relativity: the framework}.
\bibitem{JanzenBoundary} D.~Janzen, \emph{The matter boundary: what the substrate's symmetry can and cannot carry}.
\bibitem{JanzenGeometricCore} D.~Janzen, \emph{The geometric core: the substrate as a real object, and maximal symmetry worn seven ways}.
\end{thebibliography}

\clearpage
\section*{Appendix L\quad The Knowledge Ledgers}\label{app:ledgers}
\addcontentsline{toc}{section}{Appendix L: The Knowledge Ledgers}
\small\noindent Each claim marked \ldgmarker\ in the text rests on a field bake or the figure--theorem bake recorded in the named ledger, which ships with this work and carries the probes, the receipts, and the three registers kept apart --- what bit, what bounced, and where the boundary is. A ledger is working apparatus, not a publication, and is cited here rather than in the bibliography. This appendix is generated from the ledger index, not maintained by hand.\par\medskip
\begingroup\renewcommand{\arraystretch}{1.25}
\begin{description}
\item[\label{ldg:algebraic_geometry}\textbf{L2}\quad\texttt{ALGEBRAIC\_GEOMETRY\_LEDGER.md}]\hfill\textsf{[field bake]}\\
The algebraic-geometry field-bake ledger --- what bit, what bounced, and the boundary. One of the three fields listed but never thrown (the overnight order: algebraic geometry $\times$57, discriminant $\times$28 / genus $\times$21). `OWED` 622
\item[\label{ldg:functional_analysis}\textbf{L10}\quad\texttt{FUNCTIONAL\_ANALYSIS\_LEDGER.md}]\hfill\textsf{[field bake]}\\
The functional-analysis / unitarity field-bake ledger --- the field that bounced, and the one routing fact it returned. Third of the four fields `L-272`'s re-survey left outstanding. `OWED` 622
\item[\label{ldg:harmonic_analysis}\textbf{L12}\quad\texttt{HARMONIC\_ANALYSIS\_LEDGER.md}]\hfill\textsf{[field bake]}\\
The harmonic-analysis field-bake ledger --- what bit, what bounced, and the boundary. Second of the four fields `L-272`'s re-survey left outstanding. `OWED` 622
\item[\label{ldg:quadric_geometry}\textbf{L17}\quad\texttt{QUADRIC\_GEOMETRY\_LEDGER.md}]\hfill\textsf{[field bake]}\\
projective geometry of quadrics against the CR substrate --- the CK metric identification, and the ladder gap it exposes
\item[\label{ldg:representation_theory}\textbf{L18}\quad\texttt{REPRESENTATION\_THEORY\_LEDGER.md}]\hfill\textsf{[field bake]}\\
The representation-theory field bake --- what bit, what bounced, and the boundary. The largest unbaked vocabulary in the corpus ($\times$241 tight), thrown after the Phase 4 survey named it the standing first pick on measured usage
\item[\label{ldg:spectral_theory}\textbf{L19}\quad\texttt{SPECTRAL\_THEORY\_LEDGER.md}]\hfill\textsf{[field bake]}\\
The spectral-theory field bake --- what bit, what bounced, and the boundary. Tier B's largest never-thrown field ($\times$189 on the reach measure), verified as NOT covered by the harmonic ledger, which mentions spectral, self-adjoint and deficiency zero times
\item[\label{ldg:information_theory}\textbf{L21}\quad\texttt{INFORMATION\_THEORY\_LEDGER.md}]\hfill\textsf{[field bake]}\\
The information-theory field-bake ledger --- what bit, what bounced, and where the boundary is
\item[\label{ldg:number_theory}\textbf{L16}\quad\texttt{NUMBER\_THEORY\_LEDGER.md}]\hfill\textsf{[field bake]}\\
The number-theory field-bake ledger --- what bit, what bounced, and where the boundary is
\item[\label{ldg:probability}\textbf{L24}\quad\texttt{PROBABILITY\_LEDGER.md}]\hfill\textsf{[field bake]}\\
The probability and stochastic-processes field-bake ledger --- what bit, what bounced, and where the boundary is
\end{description}
\endgroup


\end{document}
